\chapter*{Índice de conexiones con la literatura de investigación}
\addcontentsline{toc}{chapter}{Índice de conexiones con la literatura de investigación}

Esta tabla permite regresar a una fuente después de estudiar el concepto del curso.
La palabra \emph{relacionada} no afirma que el objeto avanzado sea una
generalización del resultado elemental. En las filas de operadores,
$X$ denota un espacio de Banach complejo y $H$ uno de Hilbert.

\footnotesize
\begin{longtable}{>{\raggedright\arraybackslash}p{2.5cm}>{\raggedright\arraybackslash}p{5.5cm}>{\raggedright\arraybackslash}p{3.6cm}>{\raggedright\arraybackslash}p{2.0cm}}
\toprule
Concepto del curso & Aparición y qué reconocer & Fuente precisa & Tipo\\
\midrule
\endfirsthead
\toprule
Concepto del curso & Aparición y qué reconocer & Fuente precisa & Tipo\\
\midrule
\endhead
Álgebra de Banach &
$T=V_\alpha P_++M_GP_-\in\mathcal B(L^p(\Gamma))$:
suma y composición de operadores acotados. &
Karlovich 2017, pp.~463--464, (1.1)--(1.2) \cite{Karlovich2017}. &
Ejemplo directo\\
Álgebra $C^*$ conmutativa &
$C_b(\Gamma)$ y $\SO(\Gamma)$ son álgebras concretas de coeficientes,
con conjugación puntual. &
Karlovich 2017, p.~464 \cite{Karlovich2017}. &
Ejemplo directo\\
Ideal y cociente &
$J=\mathcal K(X)$ en $\A=\mathcal B(X)$ da la álgebra de Calkin;
la invertibilidad de la clase caracteriza Fredholmness. &
Karlovich 2025, p.~2 \cite{Karlovich2025}. &
Uso de propiedad\\
Espectro &
$\sigma_{\mathrm{ess}}(A)$ se define aquí mediante la clase en Calkin:
se distingue de $\sigma_{\mathcal B(X)}(A)$. &
Elaboración de estas notas a partir del cociente citado en Karlovich 2025, p.~2 \cite{Karlovich2025}. &
Construcción relacionada\\
Caracteres y Gelfand &
Las fibras $M_0,M_\infty$ de $\Omega(\SO(\mathbb R_+))$ codifican
comportamientos límite de coeficientes. &
Karlovich 2017, §2, pp.~465--466 \cite{Karlovich2017}. &
Uso de propiedad\\
Proyección en $C^*$ &
En $L^2(\Gamma)$, para $\Gamma=\mathbb R$ o $\mathbb T$,
$P_\Gamma^\pm$ son ortogonales.
La misma fórmula en la estrella abierta no garantiza idempotencia. &
BFK 2006, §2, p.~22 \cite{BFK2006}; Karlovich 2017, teorema~5.3, p.~476 \cite{Karlovich2017}. &
Ejemplo con hipótesis\\
Unitarios y shifts &
Las rotaciones $U_z$ son unitarias; el shift $V_\alpha$ de Haseman
requiere otras hipótesis de acotación. &
BFK 2006, §1, p.~20 \cite{BFK2006}; Karlovich 2017, §2, p.~467 \cite{Karlovich2017}. &
Comparación\\
Álgebra generada &
Multiplicadores, $S_\Gamma$ y conjugación generan una álgebra real
de operadores; tipar la conjugación evita llamarla compleja. &
Karlovich 2025, (1.4), p.~3 \cite{Karlovich2025}. &
Ejemplo directo\\
Homomorfismos &
Los mapas matriciales de símbolo ayudan a comprobar Fredholmness
en una álgebra concreta; no son automáticamente $*$-homomorfismos. &
Karlovich 2025, teoremas~6.2 y 7.2, pp.~25--26 \cite{Karlovich2025}. &
Construcción relacionada\\
Fredholmness &
El criterio para $T$ llega tras la representación Mellin módulo
compactos y sus hipótesis de símbolo. &
Karlovich 2017, teorema~7.1, pp.~487--489 \cite{Karlovich2017}. &
Para leer más\\
\bottomrule
\end{longtable}
\normalsize
